% Tableau des labs du cours ACSI dans NDAP SUKULU (propre au cours).
Le cours compte dix laboratoires, disponibles sous le code \codeUE. Le dernier
réunit les trois cas du fil rouge et sert de synthèse.

\begin{tableaucours}{lL{6.0cm}L{1.6cm}L{2.0cm}}{%
    \tableauHeader{Lab} &
    \tableauHeader{Sujet} &
    \tableauHeader{Chapitre} &
    \tableauHeader{Opération}
}
ACSI-LAB-01 & Construire la cartographie des risques & \ref{chap:risque} & LIZBIZ \\
ACSI-LAB-02 & Analyser la matrice des droits & \ref{chap:controle-interne} & LIZBIZ \\
ACSI-LAB-03 & Coter les documents sollicités, dimensionner les échantillons &
\ref{chap:planifier} & LIZBIZ \\
ACSI-LAB-04 & Contrôler une décomposition du domaine D & \ref{chap:decomposition} &
LIZBIZ \\
ACSI-LAB-05 & Mesurer, situer et comparer les domaines & \ref{chap:conclure} & LIZBIZ \\
ACSI-LAB-06 & Audit de conformité des paiements & \ref{chap:droit} & MEFOU \\
ACSI-LAB-07 & Tests substantifs sur les données & \ref{chap:audit-applications} &
LIZBIZ \\
ACSI-LAB-08 & Audit d'infrastructure des quatre centres & \ref{chap:audit-securite}
& LIZBIZ \\
ACSI-LAB-09 & Auditer la continuité du système central & \ref{chap:audit-continuite}
& SANAGA \\
ACSI-LAB-10 & Synthèse des trois cas & \ref{chap:cas-integrateurs} & les trois \\
\end{tableaucours}
\vspace{-0.8em}
\begin{center}{\small\itshape Les laboratoires du cours dans NDAP SUKULU}\end{center}
